\begin{frame}[t]
  \framesubtitle{Appendix}
  \frametitle{Backup roadmap}
  \begin{columns}[T,totalwidth=\textwidth]
    \column{0.48\textwidth}
    \begin{itemize}
      \item DDPM vs stochastic interpolants
      \item velocity / score relations
      \item CKNNA and linear probing
      \item DiT/SiT block details
      \item DiT analysis and more ablations
    \end{itemize}
  \column{0.48\textwidth}
    \begin{itemize}
      \item hyperparameters and compute
      \item detailed quantitative tables
      \item 512x512 extension
      \item text-to-image extension
      \item feature maps and extra qualitative samples
    \end{itemize}
  \end{columns}
\end{frame}

\begin{frame}[t]
  \framesubtitle{Appendix}
  \frametitle{DDPM vs stochastic interpolants}
  \begin{columns}[T,totalwidth=\textwidth]
    \column{0.48\textwidth}
    \textbf{DDPM view}
    \begin{itemize}
      \item discrete-time forward corruption with \(\beta_t\)
      \item learn reverse transitions \(p(x_{t-1}\mid x_t)\)
      \item common objective predicts noise \(\epsilon\)
    \end{itemize}
    \begin{align*}
      \mathcal{L}_{\text{simple}}
      = \mathbb{E}_{x^\ast,\epsilon,t}\left\|\epsilon-\epsilon_\theta(x_t,t)\right\|_2^2
    \end{align*}
  \column{0.48\textwidth}
    \textbf{Stochastic interpolant view}
    \begin{itemize}
      \item continuous path between data and Gaussian noise
      \item learn velocity or score over continuous time
      \item cleaner interface for flow-matching style models such as SiT
    \end{itemize}
    \begin{align*}
      x_t = \alpha_t x^\ast + \sigma_t \epsilon,\qquad
      \dot{x}_t = v(x_t,t)
    \end{align*}
    \vspace{0.12cm}
    {\small REPA is compatible with both views because it attaches to hidden representations rather than a specific sampler.}
  \end{columns}
\end{frame}

\begin{frame}[t]
  \framesubtitle{Appendix}
  \frametitle{Velocity and score relations}
  \begin{align*}
    v(x,t)
    &= \dot{\alpha}_t \mathbb{E}[x^\ast \mid x_t=x] + \dot{\sigma}_t \mathbb{E}[\epsilon \mid x_t=x] \\
    \mathcal{L}_{\text{velocity}}(\theta)
    &= \mathbb{E}_{x^\ast,\epsilon,t}
    \left\|v_\theta(x_t,t)-\dot{\alpha}_t x^\ast-\dot{\sigma}_t \epsilon\right\|_2^2
  \end{align*}
  \begin{align*}
    s(x,t) &= -\sigma_t^{-1}\mathbb{E}[\epsilon\mid x_t=x] \\
    s(x,t) &= \sigma_t^{-1}\cdot
    \frac{\alpha_t v(x,t)-\dot{\alpha}_t x}
         {\dot{\alpha}_t \sigma_t - \alpha_t \dot{\sigma}_t}
  \end{align*}
  \vspace{0.1cm}
  \callout{\textbf{Why these equations matter for REPA.} The paper keeps the denoising math intact.
  REPA only adds feature supervision on top of the same hidden states used to solve these objectives.}
\end{frame}

\begin{frame}[t]
  \framesubtitle{Appendix}
  \frametitle{What are CKNNA and linear probing measuring?}
  \begin{columns}[T,totalwidth=\textwidth]
    \column{0.52\textwidth}
    \textbf{Linear probing}
    \begin{itemize}
      \item Freeze the representation.
      \item Train a linear classifier on ImageNet labels.
      \item Higher accuracy implies stronger semantic separability.
    \end{itemize}
    \vspace{0.16cm}
    \textbf{CKNNA}
    \begin{itemize}
      \item A relaxed alignment metric related to CKA.
      \item Focuses on mutual nearest neighbors rather than all pairwise relationships.
      \item More forgiving when comparing representations from different model families.
    \end{itemize}
  \column{0.44\textwidth}
    \begin{align*}
      \mathrm{CKA}(K,L)
      =
      \frac{\mathrm{HSIC}(K,L)}
      {\sqrt{\mathrm{HSIC}(K,K)\mathrm{HSIC}(L,L)}}
    \end{align*}
    \begin{beamercolorbox}[wd=\linewidth,rounded=true,sep=0.18cm]{block body}
      \textbf{How the paper uses them}
      \begin{itemize}
        \item probing: semantic gap
        \item CKNNA: representational alignment
      \end{itemize}
    \end{beamercolorbox}
  \end{columns}
\end{frame}

\begin{frame}[t]
  \framesubtitle{Appendix}
  \frametitle{DiT / SiT block details}
  \begin{columns}[T,totalwidth=\textwidth]
    \column{0.42\textwidth}
    \FigNine[width=\linewidth]
  \column{0.54\textwidth}
    \begin{itemize}
      \item Patchified latent tokens go through repeated attention + MLP blocks.
      \item AdaLN-zero style modulation injects timestep and condition information.
      \item REPA taps one intermediate hidden state and sends it to a separate projection MLP.
      \item The projection head is a training-time adapter into the external encoder feature space.
    \end{itemize}
  \end{columns}
\end{frame}

\begin{frame}[t]
  \framesubtitle{Appendix}
  \frametitle{The same behavior also appears in DiT}
  \begin{columns}[T,totalwidth=\textwidth]
    \column{0.47\textwidth}
    \FigTen[width=\linewidth]
  \column{0.49\textwidth}
    \begin{itemize}
      \item DiT hidden states also show a semantic gap relative to DINOv2.
      \item Alignment is present but weak, just like in SiT.
      \item So the paper's argument is not a SiT-only artifact; it generalizes across denoising transformer families.
    \end{itemize}
  \end{columns}
\end{frame}

\begin{frame}[t]
  \framesubtitle{Appendix}
  \frametitle{More ablations: timesteps, target encoders, and \(\lambda\)}
  \begin{columns}[T,totalwidth=\textwidth]
    \column{0.58\textwidth}
    \papercrop[width=\linewidth]{10}{70bp 320bp 60bp 88bp}
  \column{0.38\textwidth}
    \begin{itemize}
      \item REPA improves the representation gap across multiple noise levels.
      \item Alignment gains are not specific to DINOv2; MAE and MoCov3 also show the trend.
      \item Performance is fairly robust once \(\lambda\) reaches about \(0.5\).
    \end{itemize}
    \vspace{0.14cm}
    \begin{tabular}{C{1.0cm}C{1.0cm}C{1.0cm}C{1.0cm}}
      \toprule
      \(\lambda\) & 0.25 & 0.5 & 1.0 \\
      \midrule
      FID & 8.6 & 7.9 & 7.8 \\
      \bottomrule
    \end{tabular}
  \end{columns}
\end{frame}

\begin{frame}[t]
  \framesubtitle{Appendix}
  \frametitle{Hyperparameters and compute details}
  \begin{columns}[T,totalwidth=\textwidth]
    \column{0.60\textwidth}
    \begin{tabular}{L{2.8cm}L{3.8cm}}
      \toprule
      Category & Setup used in the paper \\
      \midrule
      Backbone & same DiT / SiT architecture as original papers \\
      Input & ImageNet 256x256, VAE latents of size \(32\times32\times4\) \\
      Optimizer & AdamW, learning rate \(1\times 10^{-4}\), no weight decay \\
      Batch size & 256 \\
      Sampler & Euler-Maruyama, \(w_t=\sigma_t\), NFE \(=250\) \\
      Projection head & 3-layer MLP with SiLU \\
      Training hardware & 8x H100 80GB, about 5.4 step/s \\
      \bottomrule
    \end{tabular}
  \column{0.34\textwidth}
    \callout{\textbf{Engineering note.} The paper mentions that training can be sped up further by precomputing pretrained-encoder features.}
  \end{columns}
\end{frame}

\begin{frame}[t]
  \framesubtitle{Appendix}
  \frametitle{Detailed quantitative results}
  \begin{columns}[T,totalwidth=\textwidth]
    \column{0.48\textwidth}
    \begin{tabular}{L{2.1cm}C{1.0cm}C{0.9cm}C{1.0cm}}
      \toprule
      Model & Iter. & FID & Acc. \\
      \midrule
      SiT-B/2 + REPA & 400K & 24.4 & 61.2 \\
      SiT-L/2 + REPA & 400K & 10.0 & 69.4 \\
      SiT-XL/2 + REPA & 400K & 7.9 & 70.3 \\
      SiT-XL/2 + REPA & 4M & 5.9 & 74.6 \\
      \bottomrule
    \end{tabular}
    \vspace{0.16cm}
    {\small Larger models gain more from REPA, and the advantage keeps growing with longer training.}
    \column{0.48\textwidth}
    \begin{tabular}{L{2.5cm}C{0.8cm}C{0.9cm}}
      \toprule
      Setting & \(w\) & FID \\
      \midrule
      Vanilla SiT-XL/2 & 1.50 & 2.06 \\
      REPA & 1.30 & 1.80 \\
      REPA + interval [0,0.75] & 2.00 & 1.78 \\
      REPA + interval [0,0.70] & 1.80 & 1.42 \\
      \bottomrule
    \end{tabular}
    \vspace{0.16cm}
    {\small Guidance interval plus REPA produces the best reported ImageNet 256x256 result in the paper.}
  \end{columns}
\end{frame}

\begin{frame}[t]
  \framesubtitle{Appendix}
  \frametitle{Extension to ImageNet 512x512}
  \begin{columns}[T,totalwidth=\textwidth]
    \column{0.60\textwidth}
    \AppFiveTwelve[width=\linewidth]
  \column{0.36\textwidth}
    \begin{itemize}
      \item Same setup idea, larger latent input.
      \item REPA already beats vanilla SiT-XL/2 with \(>3\times\) fewer iterations.
      \item Shows the method scales beyond the 256x256 headline experiment.
    \end{itemize}
  \end{columns}
\end{frame}

\begin{frame}[t]
  \framesubtitle{Appendix}
  \frametitle{Text-to-image extension}
  \begin{columns}[T,totalwidth=\textwidth]
    \column{0.64\textwidth}
    \AppTTwoI[width=\linewidth]
  \column{0.32\textwidth}
    \begin{itemize}
      \item The paper also applies REPA to MMDiT on MS-COCO.
      \item Reported FID improves from 6.05 to 4.73 in the ODE setup, and from 5.30 to 4.14 in the SDE setup.
      \item This suggests the alignment idea remains useful even when text conditioning is present.
    \end{itemize}
  \end{columns}
\end{frame}

\begin{frame}[t]
  \framesubtitle{Appendix}
  \frametitle{Feature maps and extra qualitative samples}
  \begin{columns}[T,totalwidth=\textwidth]
    \column{0.62\textwidth}
    \AppFeature[width=\linewidth]
  \column{0.34\textwidth}
    \AppQualTop[width=\linewidth]
    \vspace{0.14cm}

    \AppQualBottom[width=\linewidth]
  \end{columns}
\end{frame}
